\chapter{Wireless Mesh Networks: Background and Architecture}
\label{ch:wmn}

\section{Definition and Characteristics}

A Wireless Mesh Network (WMN) is a communication network built from radio nodes
organised in a mesh topology, where every node is capable of relaying data on behalf of
its neighbours~\cite{akyildiz2005wmn}. Unlike traditional single-hop wireless networks
in which every client communicates through a central access point, WMN nodes forward
traffic for one another over multiple wireless hops, making the network self-organising
and self-healing.

The distinguishing characteristics of WMNs are:

\begin{description}
  \item[Multi-hop communication.] Packets may traverse several intermediate relay nodes
    to reach their destination, extending the effective coverage area far beyond the range
    of any single radio transmitter.
  \item[Self-organisation.] Nodes discover their neighbours and compute routes
    automatically, without manual configuration. New nodes joining the network are
    incorporated with minimal disruption.
  \item[Self-healing.] When a link or node fails, the network re-routes traffic around
    the failure. This property is provided by the routing protocol and is one of the
    primary motivations for using WMNs in critical-infrastructure applications.
  \item[Partial infrastructure independence.] A WMN can operate without any
    pre-existing wired infrastructure, though in practice one or more gateway nodes
    connect the mesh to the Internet or a backbone network.
  \item[Scalability tension.] While multi-hop forwarding extends coverage, larger networks
    require more complex routing, more control traffic, and more relay hops that add delay
    and consume energy. Managing this tension is the central challenge in WMN design.
\end{description}

\section{Network Architecture}

WMNs are typically classified into three architectural categories~\cite{akyildiz2005wmn,
bruno2005mesh}:

\subsection{Infrastructure Mesh}

In an infrastructure WMN, a set of dedicated mesh routers forms a backbone. These
routers are typically mains-powered, have directional antennas, and run advanced
multi-radio, multi-channel configurations to maximise backbone throughput. Client devices
(laptops, smartphones, sensors) connect to the backbone through the nearest mesh router.
This architecture is used in large-scale community broadband networks and enterprise
campus networks.

\subsection{Client Mesh}

In a client WMN, the client devices themselves participate in routing. There are no
dedicated routers; every node both generates data and forwards traffic for its peers.
This architecture suits ad~hoc deployments such as disaster-relief communications and
vehicular networks (VANETs). The main challenges are device heterogeneity and energy
constraints, since the routing burden falls on battery-powered client devices.

\subsection{Hybrid Mesh}

A hybrid WMN combines dedicated mesh routers with mesh-capable client devices. Clients
can communicate directly with each other and with nearby mesh routers, while the routers
provide backbone connectivity and Internet access. This is the most common commercial
deployment model and the architecture implicitly modelled in this thesis, where all nodes
are \texttt{ManetRouter} instances capable of both application traffic and routing.

\section{Key Performance Metrics}

Evaluating a WMN routing protocol requires measuring multiple, sometimes conflicting,
performance dimensions:

\subsection{Throughput}

Throughput is the rate at which useful application data is delivered to the destination,
typically measured in bytes per second or as a total byte count over the simulation period.
High routing overhead reduces throughput by consuming channel capacity with control messages.

\subsection{End-to-End Delay}

End-to-end delay is the time elapsed between a packet being generated at the source and
being received at the destination. It includes: queuing delay (waiting for the MAC layer),
transmission delay (radio transmission time), propagation delay (signal travel time), and
routing delay (route discovery time for reactive protocols). In IoT applications, delay
requirements vary widely: industrial control requires sub-millisecond response times, while
environmental monitoring can tolerate delays of several seconds.

\subsection{Packet Delivery Ratio (PDR)}

The PDR is the fraction of transmitted packets that are successfully received at the
destination. Drops occur due to missing routes (for reactive protocols during route
discovery, or for proactive protocols during convergence after a topology change), buffer
overflow, or MAC-layer transmission failure after exhausting all retransmission attempts.

\subsection{Energy Consumption}

Energy consumption is a critical metric in IoT applications where nodes are
battery-powered. Routing protocols that generate high volumes of control traffic cause
radios to spend more time in the transmit and receive states, which are the most
energy-intensive radio states.

\subsection{Routing Overhead}

Routing overhead is the ratio of control bytes to total bytes transmitted. A high routing
overhead indicates that a large fraction of channel capacity is consumed by protocol
maintenance rather than application data.

\section{Applications of Wireless Mesh Networks}

WMNs have been deployed across a wide range of application domains:

\paragraph{Community and municipal broadband.}
Cities in Europe, Asia, and North America have deployed WMNs to provide low-cost
Internet access in public spaces and underserved areas. Notable examples include the
Madrid city-wide WMN and municipal deployments in rural Bangladesh.

\paragraph{Industrial IoT and smart manufacturing.}
Factories use WMNs to connect sensors, actuators, and programmable logic controllers
(PLCs) without the cost and rigidity of wired Ethernet. The self-healing property is
particularly valuable in environments where equipment is rearranged frequently.

\paragraph{Smart grid and utility monitoring.}
Electric utilities deploy WMNs to connect smart meters, substation monitoring equipment,
and demand-response control systems. Many-to-one traffic patterns --- all meters reporting
to a central collector --- are characteristic of this use case, and are the pattern
modelled in this thesis.

\paragraph{Disaster relief and emergency communications.}
When earthquake or flood damage destroys fixed infrastructure, first responders can deploy
portable WMN nodes in minutes to restore communications. The infrastructure-free nature
of client-mesh architectures is essential in these scenarios.

\paragraph{Vehicular Networks (VANETs).}
Vehicles equipped with IEEE~802.11p (DSRC) or cellular-V2X radios form mobile mesh
networks to exchange safety messages, traffic information, and infotainment content.
VANETs present extreme mobility challenges for routing protocols.

\paragraph{Agricultural and environmental monitoring.}
WMNs connect soil moisture sensors, weather stations, and irrigation controllers across
large fields, enabling precision agriculture without expensive cabling.

\section{Challenges in WMN Routing}

Routing in WMNs is significantly more challenging than routing in wired networks for
several reasons:

\paragraph{Shared wireless medium.}
All nodes within radio range compete for the same channel. The IEEE~802.11 CSMA/CA
protocol manages this contention, but at the cost of latency and reduced throughput as
the number of competing nodes increases.

\paragraph{Hidden terminal problem.}
Node A can hear node B, and node C can hear node B, but A cannot hear C. If A and C
transmit simultaneously, they collide at B. The IEEE~802.11 RTS/CTS mechanism
mitigates but does not eliminate this problem.

\paragraph{Many-to-one funnelling.}
In IoT data-collection scenarios, all nodes transmit to a single gateway or sink. Traffic
density increases near the sink (the hot-spot effect), creating severe MAC-layer
congestion in the nodes closest to the sink that must relay traffic from many sources.

\paragraph{Routing loop and stale-route problems.}
Distance-vector protocols are susceptible to routing loops and can maintain stale routes
that point to nodes that have moved or become unreachable. DSDV's sequence numbers
partially address this, but convergence is not instantaneous.

\paragraph{Energy heterogeneity.}
In mixed deployments, some nodes are mains-powered and others are battery-powered.
A protocol that maximises network throughput may exhaust battery-powered nodes first,
partitioning the network.
